%!TEX root = ../template.tex
%%%%%%%%%%%%%%%%%%%%%%%%%%%%%%%%%%%%%%%%%%%%%%%%%%%%%%%%%%%%%%%%%%%
%% chapter1.tex
%% NOVA thesis document file
%%
%% Chapter with introduction
%%%%%%%%%%%%%%%%%%%%%%%%%%%%%%%%%%%%%%%%%%%%%%%%%%%%%%%%%%%%%%%%%%%

\typeout{NT FILE chapter1.tex}%

\chapter{Introduction}
\label{cha:introduction}


\section{Context}
\label{sec:context}

The study of Formal Languages and Automata Theory (FLAT) is integral to 
Computer Engineering and Computer Science programs. However, its abstract 
and mathematical nature often makes the underlying concepts and models 
difficult to master. Tools that complement FLAT by enabling practical work 
and offering visual support can mitigate this difficulty.

At FCT-UNL, the OCamlFLAT/OFLAT toolset has been developed to support this area. 
Although it already covers several key models, there remain opportunities to 
strengthen its foundations and extend its capabilities. This dissertation focuses 
on the backend: it extends \textbf{OCamlFLAT} with support for attribute grammars 
and defines the interfaces required for later integration with \textbf{OFLAT}. 
The graphical treatment (user interface and visualisation) and any empirical 
pedagogical evaluation are outside the scope of this work and are identified as 
future work.


\section{Objectives}
\label{sec:objectives}

The primary objective is to incorporate support for attribute grammars into 
\textbf{OCamlFLAT}. Attribute grammars extend context-free grammars by associating 
additional information (attributes) with productions; these attributes are essential 
for expressing semantic properties of the languages generated by the grammars. 
The work provides core representations for synthesised and inherited attributes, 
evaluation strategies, and static checks for common specification errors. It also 
exposes data models to enable future visualisation and interactive editing 
in \textbf{OFLAT}.


\section{Document Structure}
\label{sec:document_structure}

This thesis is structured into eight chapters:

\begin{itemize}
  \item \textbf{Introduction}: Presents the context and motivation, states the objectives, clarifies scope and limitations (the graphical interface and any pedagogical evaluation are out of scope), and outlines the document.

  \item \textbf{Theory of Computation}: Provides the theoretical background relevant to FLAT. It introduces the Chomsky hierarchy (Types 0--3) and discusses attribute grammars (synthesized and inherited attributes, conditions, and circularity). A brief overview of automata types is also included.

  \item \textbf{Grammar Tools}: Surveys pedagogical and practical tools related to grammars and automata, focusing on \textbf{JFLAP}, \textbf{CGM}, \textbf{ANTLR}, and \textbf{YACC}, with attention to capabilities that inform the design of OCamlFLAT.

  \item \textbf{OCamlFLAT and OFLAT}: Describes the OCamlFLAT library and the OFLAT platform at a high level. It records the current state and clarifies that this dissertation implements backend extensions in \textbf{OCamlFLAT}; graphical/interface developments in \textbf{OFLAT} are outlined only as future work.

  \item \textbf{Technology}: Details the main technologies and libraries used for the backend implementation, including programming with \textbf{OCaml}, OCaml/JavaScript interoperability (e.g., \texttt{js\_of\_ocaml}), and supporting modules or data structures relevant to future integration.

  \item \textbf{OCamlFLAT Implementation}: Presents the design and implementation of attribute grammar support in OCamlFLAT: data types, the \texttt{calcAttributes} procedure, validation and acceptance functions, cycle detection, and error handling.

  \item \textbf{Unit Testing}: Describes the test suite used to validate the implementation, covering synthesized and inherited attributes, constraints, error conditions, cycle detection, and acceptance/validation scenarios.

  \item \textbf{Evaluation and Future Work}: Summarizes the outcomes and limitations, and identifies future work, including graphical integration in OFLAT and an empirical pedagogical study.
\end{itemize}



